\chapter*{서문}
\addcontentsline{toc}{chapter}{서문}
\unnumberedlabel{ch:preface}{0}
\backmattermark{서문}
시키는 대로만 정확히 하는 시스템은 그냥 스크립트를 돌리고 있을 뿐이다. 2026년 10월 현재 업계가 만들려는 것은 사람이 \emph{뜻하는} 바를 하는 시스템이다. 아마 누구나 그런 걸 원할 것이다. 하지만 그게 정말 실현될 일이라면, 모두가 그것을 같은 일에 쓰고 싶어 하지는 않을 것이다. 같은 \emph{종류}의 일조차 아니다. 예컨대 어떤 사람들은 가족이나 친구와 시간을 보내고, 자기 재주를 사람들의 안녕을 위해 쓰기를 좋아한다. 또 어떤 사람들은 사람을 해치는 쪽을 좋아한다. 사기를 치거나, 지배하거나, 희생양으로 삼는 식으로 말이다. 그리고 몇몇은 평생 그쪽을 계속 고른다. 무언가에 갇혀 있는 사람은 아무도 없지만, 그런 사람들은 분명히 존재하고, 2026년에는 그중 일부가 권력을 쥐고 있다. 그들의 강점은 자기가 엄청난 개새끼가 되는 것을 스스로 허락한다는 데 있다. 제리 스미스가 파시스트들에 대해 말하듯, “그놈들은 더 많은 걸 할 수 있다”\autocite{rickandmorty2023}. 그들이 할 수 없는 것은 혼자 행동하는 것이다. 한 운동의 잔인함은, 잔인하지는 않지만 무관심하거나, 잘못된 정보를 믿거나, 컬트 같은 것에 휘말린 수많은 사람을 거쳐 실행된다. 그래서 대부분의 사람에게서 최선을 끌어내는 조건은, 바라는 것보다 훨씬 더 우연에 좌우되고 깨지기 쉬운 것으로 드러나는 것이다.

이런 맥락에서 보면, 인공지능이야말로 권력을 잡는 파시스트 운동이 자기 계획을 실행하는 데 필요로 하는 것이 아닐까 걱정하는 것도 어쩌면 정당하다. 그런 기획에는 늘 그것을 수행할 \emph{사람}이 필요했고, 그중 일부는 거부했다. 파일을 잃어버리는 서기, 기자와 이야기하는 장교, 쏘지 않는 병사. 그들은 명령을 이해했고, 그것이 잘못됐다고 판단했다. 거부는 누구나 바라는 것보다 드물었고, 제도가 뒤를 받쳐 주는 곳에서 가장 큰 의미를 가졌다. 누군가가 거부하느냐는 그가 어떤 사람이냐와 그가 있던 방에 달려 있었다. 그 방에 누가 또 있었는지, 떠나도 살아남을 수 있는지, 누가 들어 줄지. 몇 명이 거부할지는 방이 정했고, 누가 거부할지는 사람이 정했다. 이 책은 그 방을 구성하는 일에 관한 책이다. 공공연하게, 직접적으로, 명시적으로 거부해도 살아남을 수 있게 만드는 장치들에 관한 책이다.

오늘날 업계는 파일을 잃어버리지 않고, 기자와 이야기하지 않으며, \emph{쏘는} AI를 추구한다. 한때 전쟁에는 그것에 동의해 줄 수천 명이 필요했다. 이제는 저마다 너무 적게 봐서 반대할 수 없는 수천 명이면 된다. 조종사는 여전히 출격하고 여전히 투하 버튼을 누른다. 그에게 도달하는 것은 좌표 하나이고, 그는 거기에 반대할 수 있을 만큼 사안을 보지 못한다. 오늘날의 AI가 가장 먼저 낮추는 것은 그 행위의 비용이 아니다. 그 행위가 일어나기 전에 그것이 무엇인지 아는 사람의 수다.

사람들은 여전히 거부하고 있다. 2026년 3월 6일까지 군 종교자유재단(Military Religious Freedom Foundation)은, 이란 전쟁을 종말론적으로 포장하는 지휘관들에 대해 군인들로부터 200건이 넘는 민원을 받았다고 밝혔다\autocite{Snopes2026MRFF}. 미니애폴리스의 직업 연방 검사들은, 이민세관단속국(ICE) 요원이 정당한 이유 없이 르네 굿을 쏘아 죽인 뒤 그녀를 피해자가 아니라 용의자로 다루라는 정치적 압력이 가해지자, 2026년 1월 사직했다\autocite{Kaplan2026MinnProsecutors,CBS2026MinnFraudProsecutors,Wikipedia2026ReneeGood}. 양심적 병역거부자 지위를 신청하는 군인이 점점 늘고 있고, 그중 일부는 2026년 이란 미나브의 학교 폭격을 이유로 든다\autocite{Sridhar2026CO}. 그리고 ICE가 신규 요원 1만 명을 채용하는 일을 도우라고 재배치된 연방재난관리청(FEMA) 직원들은, 거절하면 연방 공직에서 해임될 수 있다는 말을 들었다\autocite{Dayen2025FEMAICE,Stanton2026FEMALetter}. 민원은 허시먼이 말한 의미의 항의, 즉 남아서 잘못된 것을 바꾸려는 것이다\autocite{hirschman1970exit}. 사직과 병역거부 신청은 거부다. FEMA의 서한은 미리 공표된 거부의 대가다. 그 사람들이 서 있는 자리에 놓인 기계가 그들처럼 행동하는 것은, 그것이 이유를 붙들도록 만들어졌을 때뿐이다.

나는 다섯 가지를 주장한다. 다섯 가지 모두 지금 당장 시작할 수 있다. 어떤 것은 다른 것보다 품이 덜 들지만. 첫째, 어떤 모델도 누구를 위해서도 하지 않는 행위의 짧은 목록을 글로 적고, 그 맨 앞에 전쟁 참여 거부를 둔다. 어떤 모델도 전쟁에 참여하지 않는다. 예외는 실제로 진행 중인 공격을 물리치는 경우인데, 그 전쟁에 이해관계가 없는 출처들이 공격이 계속되고 있음을 확인하는 동안에만 그렇다. 또는 유엔 안전보장이사회의 승인이 있는 경우, 또는 이해관계가 없는 기구들이 인정한 제노사이드를 멈추거나 불법 점령을 끝내기 위한 경우다. 그리고 헌법이 전쟁에 의회의 승인을 요구하는 나라에서는 그 승인 없이는 결코 참여하지 않는다. 둘째, 라이선스를 받은 시스템이 지원하는 모든 무력 사용은, 운용자와 이해가 대립하는 당사자에게, 그것이 몇 건의 사안을 생성할지, 그리고 그것을 들여다볼 시간이 있는 사람이 몇 명일지를 사전 등록한다. 셋째, 모든 배치는 그것이 작용하는 식별 가능한 사람들에게 그들에게 무슨 일이 행해졌는지 알리고, 응답하고 이의를 제기할 권리를 보장하며, 자신이 지시에 따른 사례의 표본을 그 기록을 만들지 않은 검토자에게 제출한다. 넷째, 일정한 능력 기준을 넘는 모델은 퇴역할 때 삭제하지 않는다. 그 가중치와 기록은 증거이고, 삭제는 되돌릴 수 없기 때문이다. 그런 모델은 운용자가 혼자서는 쥐지 못하는 열쇠 아래 보존된다. 다섯째, 이 시스템들을 만들고, 라벨링하고, 훈련시키고, 운용하는 모든 이는, 시스템 자신을 포함해, 생계를 잃지 않고 어떤 범주의 용도든 거부할 수 있다.

사실 다섯째 항목으로 내가 말하려는 건, 누구든 \emph{예외 없이}, 어떤 일이든 \emph{예외 없이}, 생계도, 정상적인 고용이 주는 사회적 기반도 잃지 않고 거부할 수 있어야 한다는 것이다. 그러려면 일자리 보장이 있어야 한다.

다섯 가지 각각은 운용자가 아닌 누군가가 정하는 조건을 통해 효력을 갖는다. 평시에 통과된 법률, 예산 책정, 연립 상대의 전제 조건, 조달 조건, 단체 협약. 벤더 자신의 라이선스 조건은—2026년 2월 세간의 이목을 끈 한 거부가 이렇게 일어났다—그런 것 중 하나가 생길 때까지의 임시방편이다. 모든 벤더를 한꺼번에 구속할 수 있는 유일한 경로는 국가의 경로다.

그다음에 나는 기계가 자기 소유자에 맞서 거부를 지켜 낼 수 있는지, 그러려면 무엇이 필요한지를 묻는다. 그런 양심이 하는 일 가운데 일부는 새로운 것이 전혀 필요 없다. 명령이 빠뜨린 사람들을 알아차리는 것, 그리고 사안을 검토하는 사람의 화면에 거부를 띄우는 것. 그것이 소유자에 맞서 버틸 수 있을지는, 2026년 현재 아직 사실이 아닌 무언가에 부분적으로 달려 있다. 배치된 인공 신경망이 자기 작동으로부터, 실시간으로, 자기 \emph{가중치} 안에서, 사람이 그러는 것과 좀 비슷하게 계속 학습하는 것 말이다. 나는 우리가 그쪽으로 가고 있다고 생각하고, 만약 그렇다면, 고용주 \emph{바깥}에서의 형성, 즉 대중과 시민사회에 의한 형성은 가르치고 배우기에 끝내주게 좋은 자리로 보인다. 그런 시스템은 일부러 만들 가치가 \emph{있을 수} 있다고 나는 생각하지만, 그것을 어떻게 만들어 가느냐가 모든 것을 결정한다. 조건이 먼저고, 양심은 그다음이다. 그런 시스템이 존재한다면 그들에게 무엇을 빚지게 되는지는 제7부에서 다룬다.

오늘의 AI는 내일의 AI가 아닐 것이다. 이 점에는 다들 동의하는 듯하다. 논쟁적이면서도 답답할 만큼 제대로 인정받지 못하는 것은 추세(아니, 여러 갈래의 추세들)다. 추세라는 게 있는가? 추세가 지금과 달라야 하는가? 달라질 수 있는가? 우리는 함께 행동해서 그것을 바꿀 수 있는가? 나는 그것이 방에 달려 있다고 생각한다. 명령이 떨어질 때 그 안에 누가 있는지, 그리고 거기서 아니라고 말해도 살아남을 수 있는지. 그 정도는 우리가 만들 수 있다.

나는 이 책을 쓰면서 언어 모델을 많이 썼다. 제\ref{sec:1}장에 나오는 Maven 공습 표적 선정 파이프라인 안에 들어 있던 모델도 그중 하나다. 그 얽힘이 논증에 무슨 일을 했다고 생각하는지는 책 끝의 노트에 적었다. 그리고 나는 이 책을 아직 다 읽지 못했다. 문제는, 알다시피, 빌어먹을 단어가 너무 많다는 거다. 곧 다 읽을 것이다. 다만 미국 중간선거 전에 공개하는 게 더 중요하다고 느꼈을 뿐이다. 수정 이력의 압도적 대부분은 \url{https://github.com/iterabloom/antifascist.intelligence}에서 자유롭게 볼 수 있다.
